% =============================================================================
% Paper 2: The Failed Attempt to Make Transformers Continuous
% Target: NeurIPS "What Never Worked" Workshop / ICLR Meta-Research
% Format: NeurIPS-style (~30 pages main + appendix)
% =============================================================================
\documentclass{article}

\usepackage[preprint]{neurips_2024}
\usepackage[utf8]{inputenc}
\usepackage[T1]{fontenc}
\usepackage{hyperref}
\usepackage{url}
\usepackage{booktabs}
\usepackage{amsfonts}
\usepackage{nicefrac}
\usepackage{microtype}
\usepackage{xcolor}
\usepackage{amsmath}
\usepackage{graphicx}
\usepackage{algorithm}
\usepackage{algorithmic}
\usepackage{listings}
\usepackage{amsthm}

% Theorem environments
\newtheorem{theorem}{Theorem}
\newtheorem{lemma}{Lemma}
\newtheorem{corollary}{Corollary}
\newtheorem{definition}{Definition}
\newtheorem{conjecture}{Conjecture}

\title{The Failed Attempt to Make Transformers Continuous:\\
       A Systematic Record of 35 Experiments}

\author{%
  Yiming Wang \\
  CrystaLLM Project \\
  \texttt{yiming.wang@crystallm.org} \\
}

\begin{document}

\maketitle

\begin{abstract}
We document a systematic, 35-experiment campaign to test whether continuous representations can be integrated into the autoregressive Transformer. Across two research stages (CMT: 30 experiments patching KAN FFNs, complex-valued attention, and Lie group position encodings into a standard Transformer; CWF: 5 experiments designing a purpose-built closed continuous architecture), 25 experiments failed with a consistent pattern: training-inference distribution drift. We formalize this failure pattern with a \emph{closure argument}: any architecture that has both a discrete token input and a continuous internal representation cannot be fully closed on its continuous state space, and the resulting distribution mismatch causes failure. The single surviving intervention is Soft-Exp inference---a one-line code change that requires no retraining---which we report in a companion paper. We argue that the closure property is the right theoretical lens for this problem: it explains why 30+ rounds of patching failed, why Soft-Exp works, and why a future fully-continuous architecture would require 1--2 person-years of engineering investment. We publish this systematic record to save the community from repeating our experiments and to provide a foundation for the CWF research charter.
\end{abstract}

% =============================================================================
\section{Introduction}
% =============================================================================

\subsection{Motivation}

The dominant paradigm in language modeling is the autoregressive Transformer: predict the next token given previous tokens, train with teacher forcing, and deploy by feeding back the model's own predictions. This paradigm has achieved remarkable empirical success---GPT-class models, LLaMA, and many others have demonstrated that scaling autoregressive Transformers yields increasingly capable language models.

Yet a structural concern persists: the autoregressive Transformer operates on \emph{discrete} tokens at the input and output, while the underlying signal (human language) has \emph{continuous} structure (phonetic, semantic, syntactic gradients). This discretization at the interface raises a question that the field has explored sporadically but never resolved: \emph{can we make Transformers continuous?}

The intuition is appealing. Continuous representations could preserve the rich gradient information that tokenization discards, enable smooth interpolation, allow multi-modal signals to be processed without first quantizing to discrete units, and eliminate the train-inference distribution mismatch that causes exposure bias.

Over the past five years, several research communities have pursued this intuition through diverse approaches: complex-valued neural networks, Kolmogorov-Arnold representations, Lie group actions, neural ODEs, and diffusion language models. Each has had promising theoretical motivations. Each has, in our experience, failed to deliver on the autoregressive Transformer.

\subsection{Our Method}

We conducted two research campaigns:

\textbf{CMT Campaign (Exp 1--30, $\sim$3 months)}: patch continuous components onto a standard autoregressive Transformer. We tested 10+ combinations of KAN FFNs, complex-valued attention, and Lie group position encodings across 5 model scales (2M to 1.2B parameters) and multiple tokenization schemes (character-level and BPE). The campaign concluded with the \textbf{CMT\_DEAD\_FINAL} verdict after 30 rounds.

\textbf{CWF Campaign (Exp 01--02, 2 weeks)}: design a fully closed continuous architecture from scratch. We formulated a closure conjecture, implemented a minimal CWF prototype with closure verified across training, and stress-tested it on the Lorenz chaotic system. The campaign concluded with a \textbf{TOTAL\_FAIL} at the GO/NO-GO gate.

Both campaigns were data-driven: each experiment's design was informed by the previous experiment's results.

\subsection{The 35 Experiments at a Glance}

\begin{table}[h]
\centering
\small
\begin{tabular}{lcc}
\toprule
Verdict & Count & Percentage \\
\midrule
STRONG PASS & 1 & 2.9\% \\
PASS & 4 & 11.4\% \\
PARTIAL & 1 & 2.9\% \\
INCONCLUSIVE & 4 & 11.4\% \\
\textbf{FAIL} & \textbf{25} & \textbf{71.4\%} \\
\bottomrule
\end{tabular}
\caption{Distribution of verdicts across the 35 experiments.}
\label{tab:verdicts}
\end{table}

The single STRONG PASS is Soft-Exp on V49 1.2B \citep{wang2026softexp}. The 25 FAILs span the entire design space we explored.

\subsection{The Three Independent Evidence Streams}

We arrived at the "patching is futile" conclusion from three independent streams:

\textbf{Stream 1: CMT Campaign (30 experiments).} We patched KAN FFNs, complex-valued attention, and Lie group position encodings into autoregressive Transformers in 10+ combinations. All failed with a consistent pattern: training-inference distribution drift, expressed as memorization at 8--12k training steps.

\textbf{Stream 2: Soft-Exp Training Test (Exp 30).} We took the Soft-Exp inference method, which works ($+48.6\%$), and applied it at training time. The result was catastrophic. This is the cleanest possible demonstration that the training-inference gap must be closed from the inference side.

\textbf{Stream 3: CWF Campaign (5 experiments).} We abandoned patching and designed a purpose-built closed continuous architecture. Despite closure being verified, the architecture failed to learn the Lorenz dynamics. The root cause was the absence of ODE evolution machinery.

\subsection{The Closure Argument (Preview)}

We define \emph{closure} as the property that every layer maps the representation space into itself. For continuous-internal models, we propose the open unit ball $\mathbb{D}^d$ as the representation space.

\begin{theorem}[Patchwork continuity is impossible]
No architecture that has both (a) a discrete token input and (b) a continuous internal representation can be fully closed on $\mathbb{D}^d$.
\end{theorem}

The proof rests on the observation that a discrete token embedding maps a finite vocabulary to a finite subset of $\mathbb{D}^d$, while a continuous internal representation requires the entire $\mathbb{D}^d$. The corollary explains Soft-Exp's success: Soft-Exp's expected embedding lies in the \emph{convex hull} of the discrete embedding image.

\subsection{Contributions}

\begin{enumerate}
    \item \textbf{Systematic empirical record}: 35 experiments across two campaigns, documented with full reproducibility.
    \item \textbf{Five cross-cutting failure modes}: identified from the experiments, each with formal definition.
    \item \textbf{Closure theorem}: formal argument for why patchwork cannot be closed.
    \item \textbf{Lorenz decisive experiment}: clean stress test with three models.
    \item \textbf{Research charter}: the CWF Manifesto as an open problem for the community.
\end{enumerate}

\subsection{Outline}

\S2 reviews the theoretical motivation. \S3 narrates the CMT Campaign. \S4 introduces the CWF Pivot. \S5 presents the closure argument. \S6 reports the CWF Phase 1 experiment. \S7 discusses the Soft-Exp survivor. \S8 concludes. Appendices A--C provide failure modes, reproducibility, and timeline.

% =============================================================================
\section{The Continuous Wave Function Vision}
% =============================================================================

\subsection{Why continuous representations?}

Three theoretical intuitions motivate continuous representations:

\textbf{(1) Information preservation.} Discrete tokenization discards gradient information. The cosine similarity between ``run'' and ``slowly running'' is not represented in the discrete vocabulary, even though both share semantic structure.

\textbf{(2) Multi-modal integration.} Continuous signals (audio waveforms, video frames, sensor streams) cannot be processed by discrete token models without quantization, which loses information.

\textbf{(3) Geometric closure.} Continuous representations live on a manifold, enabling Lie group actions (rotations, translations) that preserve distances. Discrete tokens have no natural manifold structure.

\subsection{Theoretical foundations}

The vision of continuous language models draws on several mathematical traditions:

\textbf{Complex-valued neural networks} \citep{nitta2003complex,trabelsi2018deep}: complex-valued weights represent magnitude and phase separately, doubling the effective channel capacity per parameter.

\textbf{Kolmogorov-Arnold Networks (KAN)} \citep{liu2024kan}: learn continuous activation functions on edges rather than nodes, providing parameter efficiency for smooth functions.

\textbf{Lie group representations} \citep{lezama2018overcoming,wang2023rope}: use skew-symmetric matrices to generate rotations that preserve norms, enabling continuous position encoding.

\textbf{Born rule} \citep{born1926quantenmechanik}: probability amplitude is complex-valued; measurement gives real-valued probability via $|\psi|^2$. This suggests a quantum-mechanical analogy for language: the internal state is complex, the output is the Born-rule measurement.

\subsection{The CMT intuition}

The first wave of continuous-Transformer research (the ``CMT intuition'') proposed to \emph{patch} continuous components into a standard autoregressive Transformer:

\begin{itemize}
    \item \textbf{Knife 1}: replace the GELU MLP with a complex-valued KAN.
    \item \textbf{Knife 2}: replace softmax attention with complex-valued interference.
    \item \textbf{Knife 3}: replace RoPE with Lie group rotations.
    \item \textbf{Knife 4}: replace argmax+embedding feedback with continuous expected embedding at inference.
\end{itemize}

The hope: each knife adds continuous structure without disrupting the autoregressive training pipeline. The reality, as the next 30 experiments show, was different.

% =============================================================================
\section{The CMT Campaign (30 Experiments)}
% =============================================================================

\subsection{Pre-CMT Explorations (Exp 1--5)}

Before committing to the CMT architecture, we conducted five exploratory experiments to scope the design space.

\textbf{Exp 1 (Mamba3/SSD)}: tested whether Mamba-style selective state-space models could replace attention. The hypothesis was that SSMs' continuous-time dynamics might avoid discrete-tokenization issues. Result: perplexity higher than baseline; abandoned.

\textbf{Exp 2 (Complex KAN)}: tested KAN with complex-valued weights as an alternative to GELU MLPs. Result: strictly worse than GELU MLP.

\textbf{Exp 3 (FP8 mixed precision)}: tested FP8 forward, FP32 backward. Orthogonal to continuous-vs-discrete; ran as sanity check.

\textbf{Exp 4 (8-bit AdamW)}: tested 8-bit optimizer states. Orthogonal; ran to enable larger model experiments.

\textbf{Exp 5 (Curriculum learning)}: tested easy-to-hard sample ordering. No improvement.

Lesson: do not pursue Mamba-style alternatives or KAN-style activations independently.

\subsection{CMT Main Campaign (Exp 6--15)}

The first round combined three ``knives'':
\begin{itemize}
    \item Knife 1: KAN-based complex-valued FFN
    \item Knife 2: Complex-valued attention
    \item Knife 3: Lie group rotation for position encoding
\end{itemize}

\textbf{Exp 6 (CMT-FFN only)}: replaced only the FFN. PPL rose from baseline 2.8 to 4.5. First indication that complex-valued components were a liability.

\textbf{Exp 7 (CMT-full sanity)}: combined all three knives. At 4k steps, the model collapsed to memorizer (PPL $< 1.0$). Suspected a bug; debugging began.

\textbf{Exp 8 (CMT-full refined)}: fixed bugs in complex multiplication and Lie algebra parameterization. Still collapsed to memorizer by 8k steps.

\textbf{Exp 9 (Baseline rerun)}: confirmed baseline stable. Failure specific to CMT.

\textbf{Exp 10--13} tested surgical fixes: complex-softmax approximation, complex multiplication correctness, Lie algebra initialization, removing context-aware PE. Each gave marginal improvement; model still collapsed at 12--15k steps.

\textbf{Exp 14 (no context PE)}: removing context-aware PE gave brief improvement (PPL 32.58 peak) but model still fell to memorization by 8k.

\textbf{Exp 15 (CMT-full v2)}: combining all fixes, peak PPL was 1.01, identical to memorization.

\textbf{Diagnostic insight}: Exp 14's peak PPL of 32.58 was the highest the CMT architecture ever achieved. After 8k steps, PPL would drop to memorization values. This was a \emph{phase transition}: the model learns the training distribution's ``average behavior'' first, then sharpens into memorization.

\subsection{CMT Engineering Audit (Exp 16--25)}

Faced with the persistent memorization pattern, we launched a systematic audit.

\textbf{Exp 16 (CMT-clean)}: re-implemented CMT from scratch with all known bugs fixed. After 30k steps, val PPL reached 1.0097 (memorizer). Even a ``bug-free'' CMT failed. Turning point: failure no longer attributable to implementation errors.

\textbf{Exp 17 (phase transition diagnosis)}: trained three CMT variants at different learning rates, evaluating at 1k, 2k, 4k, 8k. Two variants (A0, A2) collapsed to memorization by 4k. The third (A1, LR $3 \times 10^{-5}$) showed partial learning but was still underfit at 8k. Phase transition was real and LR-driven.

\textbf{Exp 18 (A1 long training)}: extended A1 to 8k. Four checkpoints all underfit (val\_ppl 17.4 $\to$ 11.8, still decreasing). A1 was not converged.

\textbf{Exp 19 (5k sanity)}: small-data test (2M, 2k samples). Val PPL dropped from 33 to 18, but repetition 3/6, zero coherence. Dual-phase learning curve confirmed.

\textbf{Exp 20 (BPE sanity)}: switching to BPE tokenization gave the first LM signal in the CMT campaign: coherence 0/6 $\to$ 2/6, bits/char 4.17 $\to$ 3.32. BPE helped somewhat.

\textbf{Exp 21 (Baseline + BPE 5k)}: non-CMT baseline achieved coherence 1/6 at 5k. CMT-BPE comparison inconclusive.

\textbf{Exp 22 (CMT + BPE 16k)}: extending CMT+BPE to 16k gave a surprising result: val PPL \emph{rebounded} from 786 to 3860, while coherence rose from 0 to 5/6. We initially interpreted as ``structural learning''.

\textbf{Exp 23 (CMT + BPE 10k + 16k)}: re-ran with $5\times$ more data, $5\times$ fewer epochs. Rebound disappeared (ratio $3.90 \to 0.00$), val PPL dropped to 126. \textbf{Exp 22 was a memorization illusion}: 64 epochs over 2k samples caused sequence memorization.

\textbf{Exp 24 (Cayley PE isolation)}: isolated Cayley PE from the rest of CMT, comparing to RoPE and NoPE. Cayley PE was not better than RoPE, but not worse either. Partial success: Cayley PE alone did not break the model.

\textbf{Exp 25 (CMT + BPE + 10k + 32k)}: final CMT test. CMT reached PPL 1.40 at 12k (memorizer); baseline reached PPL 50.92 at 32k (real LM). \textbf{CMT\_DEAD\_FINAL}.

\subsection{The Soft-Exp Survivor (Exp 26--29)}

The pivot was not architectural but philosophical: instead of asking ``how do we add continuous components to Transformers?'', we asked ``can we use continuous \emph{feedback} signals without changing the architecture?''

The answer was Soft-Exp: at inference time, replace the argmax token embedding with the continuous probability-weighted sum of all token embeddings.

\textbf{Exp 26 (2M)}: Soft-Exp yielded $+81\%$ relative improvement over argmax.

\textbf{Exp 27 (8M)}: improvement was $+53\%$.

\textbf{Exp 28 (16M/32M)}: at 16M with fixed LR schedule, $+48.6\%$. At 32M (underfit), $+11\%$.

\textbf{Exp 29 (V49 1.2B)}: improvement was $+48.6\%$, identical to 16M. Scale-invariant across $600\times$ parameter range.

\subsection{The Final Test (Exp 30): Soft-Exp at Training}

If Soft-Exp worked at inference, what about training? We tested: replace ground-truth embedding with continuous expected embedding, linearly warmed from 0 to 0.5 over 1000 steps.

Result: \textbf{catastrophic failure}. Validation PPL rose to 8554 on V49 1.2B. The argmax-feedback PPL rose to 10662. The model was damaged in its core language modeling capacity.

This was the decisive experiment. Soft-Exp's success was specifically because it did not modify the training distribution. When the training input distribution was modified to include continuous expected embeddings, the model overfit to a distribution that was OOD with respect to any single inference distribution.

\subsection{Lessons from the CMT Campaign}

\begin{enumerate}
    \item Patching continuous components into autoregressive Transformers does not work.
    \item Bug-fixing does not fix structural problems.
    \item Continuous feedback at inference works; continuous feedback at training does not.
    \item The Soft-Exp one-line fix outperforms 30 rounds of architectural engineering.
\end{enumerate}

These lessons motivated the second campaign: design a continuous architecture from scratch.

% =============================================================================
\section{The CWF Pivot: From Patching to Redesign}
% =============================================================================

\subsection{Why Patchwork Was a Dead End}

The CMT Campaign established that no combination of KAN FFNs, complex-valued attention, and Lie group position encodings could be added to a standard autoregressive Transformer without catastrophic failure. We documented five distinct failure modes (Appendix A), each reproducible across multiple experiments:

\begin{enumerate}
    \item Training-inference distribution drift
    \item Closure violation in patchwork architectures
    \item Architectural absence of ODE evolution
    \item Tokenization bottleneck on continuous data
    \item Optimizer-landscape mismatch on complex losses
\end{enumerate}

Each failure mode is structural, not implementation-specific. After 30 rounds of fixes, we concluded: \emph{the patching approach is exhausted}.

\subsection{The Pivot Decision}

In May 2026, we made the decision to abandon patching and design a continuous architecture from scratch. The reasoning:

\begin{itemize}
    \item The closure property had emerged as the key theoretical concept.
    \item No existing architecture is closed on a continuous representation space.
    \item A purpose-built closed architecture might avoid the CMT failure modes by construction.
\end{itemize}

We formalized this as the \textbf{Closure Conjecture} and designed the \textbf{Closed Waveformer (CWF)} as a candidate architecture.

\subsection{The Closure Conjecture}

\begin{conjecture}[CWF Closure]
Let $\psi \in \mathbb{D}^d$ be a wave function state. Let $\mathcal{F}_l: \mathbb{D}^d \to \mathbb{D}^d$ be any of the six component operators. Then:
\begin{enumerate}
    \item $\mathcal{F}_l$ is well-defined (maps into the open disk).
    \item $\mathcal{F}_l$ is smooth.
    \item The composition $\mathcal{F}_N \circ \cdots \circ \mathcal{F}_1$ preserves $\|\psi\| < 1$.
    \item The gradient $\nabla_\theta \mathcal{L}$ exists and is bounded.
\end{enumerate}
\end{conjecture}

\subsection{The Six-Component Architecture}

We designed CWF as a stack of six closed components.

\textbf{Component 1: Continuous Input Encoding}. $\mathcal{E}(x) = \mathrm{FFT}(x) / \|x\|_2$. Frequency-domain embedding, natural for periodic signals.

\textbf{Component 2: Lie Group Position Rotation}. $\mathcal{P}_\theta(\psi) = \exp(A(\psi)) \cdot \psi$, with $A \in \mathfrak{so}(d)$. Cayley transform for efficiency.

\textbf{Component 3: Complex-Valued Attention}. $\mathcal{A}_\theta(\psi)_i = \sum_j \frac{\langle \phi_i, \psi_j \rangle}{|\langle \phi_i, \psi_j \rangle|} \cdot V_j$. Unit-phase weights preserve phase information.

\textbf{Component 4: Complex KAN-FFN}. $\mathcal{K}_\theta(\psi) = \sum_k \alpha_k \cdot B_k(\psi)$. Spectral norm constraints maintain bounded output.

\textbf{Component 5: Born-Stable Normalization}. $\mathcal{N}(\psi) = \psi / \max(\|\psi\|, 1 - \epsilon)$. Hard projection to $\mathbb{D}^d$.

\textbf{Component 6: Born Rule Output}. $P(v) = |\langle \Phi_v, \psi \rangle|^2 / \sum_v |\langle \Phi_v, \psi \rangle|^2$. Decoding only at the final layer.

\subsection{The Prototype (Phase 0)}

We implemented a minimal CWF prototype. Closure was verified across 10 random batches at initialization. The prototype learned the harmonic sequence task: training loss decreased from 320 to 3 over 2000 steps. It was 65$\times$ slower than AR baseline. Relative loss was $2.5\times$ AR baseline at 500 steps, but CWF converged faster (90 vs 190 steps).

Phase 0 demonstrated that closure is achievable. It did not demonstrate that closure is sufficient for chaotic dynamics.

% =============================================================================
\section{Why "Patching" Fails: A Mathematical Argument}
% =============================================================================

\subsection{The Closure Property}

\begin{definition}[Closure]
A neural network layer $f_\theta: \mathcal{X} \to \mathcal{X}$ is \emph{closed} on $\mathcal{X}$ if $\forall x \in \mathcal{X}, f_\theta(x) \in \mathcal{X}$.
\end{definition}

For standard Transformers, $\mathcal{X} = \mathbb{R}^d$ and closure is automatic. For continuous-internal models, we propose $\mathcal{X} = \mathbb{D}^d$.

\begin{definition}[Full closure]
A neural network is \emph{fully closed} on $\mathbb{D}^d$ if every layer maps $\mathbb{D}^d \to \mathbb{D}^d$.
\end{definition}

\subsection{Why Patchwork Violates Closure}

\begin{lemma}[Discrete-continuous interface breaks closure]
Let $E: \{0, 1, \ldots, V-1\} \to \mathbb{D}^d$ be a token embedding. If $E$ is bijective onto a finite subset of $\mathbb{D}^d$, then for any $z \in \mathbb{D}^d$ not in the image of $E$, no layer downstream of $E$ can recover the discrete structure without leaving $\mathbb{D}^d$.
\end{lemma}

\begin{proof}
$E$ maps the discrete vocabulary to a finite set $\{e_0, \ldots, e_{V-1}\} \subset \mathbb{D}^d$. Any $z \in \mathbb{D}^d$ not in this set has no discrete preimage. Downstream layers cannot determine which discrete token $z$ corresponds to without a decoding function. The decoding function $D: \mathbb{D}^d \to \mathbb{R}^V$ is non-injective. Therefore, the composed map $D \circ E$ loses information about $z$ outside the finite image. \qedhere
\end{proof}

\begin{theorem}[Patchwork continuity is impossible]
No architecture with both (a) a discrete token input and (b) a continuous internal representation can be fully closed on $\mathbb{D}^d$.
\end{theorem}

\begin{proof}
By Lemma 5.2, the input boundary $E$ maps discrete tokens to a finite subset of $\mathbb{D}^d$. For closure to hold, all internal layers must map this finite subset to $\mathbb{D}^d$ (achievable). However, the \emph{converse} closure property fails: not every $z \in \mathbb{D}^d$ has a discrete preimage. Internal layers cannot guarantee outputs lie in the finite image of $E$, and the input layer cannot recover discrete structure from arbitrary $z$. \qedhere
\end{proof}

\begin{corollary}[Training-inference drift]
The model's training input distribution is supported on $\{e_0, \ldots, e_{V-1}\}$ (discrete embedding image), while the inference input distribution is supported on a broader subset of $\mathbb{D}^d$. The two distributions are not identical; inference is \emph{broader} than training.
\end{corollary}

This is the formal statement of the training-inference distribution drift observed in Exp 30 and the CMT memorization pattern.

\subsection{Why Soft-Exp Avoids Drift}

\begin{lemma}[Soft-Exp preserves discrete support]
The continuous expected embedding $\mathbf{e}_{\text{exp}} = \mathbf{p}^\top \mathbf{E}$ lies in the convex hull of $\{e_0, \ldots, e_{V-1}\}$.
\end{lemma}

\begin{proof}
$\mathbf{e}_{\text{exp}}$ is a convex combination (probabilities non-negative, sum to 1) of the $e_v$'s, hence in their convex hull. \qedhere
\end{proof}

The convex hull of the discrete embedding image is a subset of $\mathbb{D}^d$. Soft-Exp constrains the inference input distribution to this convex hull, which is a subset of (not broader than) the training distribution's support.

\begin{theorem}[Inference narrowing is safe]
If the inference input distribution's support is a subset of the training input distribution's support, the model is in-distribution at inference time.
\end{theorem}

\subsection{The Closure Verdict}

Combining the above:

\begin{itemize}
    \item Patchwork architectures violate closure (Theorem 5.3).
    \item Continuous internal representations cause training-inference drift (Corollary 5.4).
    \item Drift causes memorization or catastrophic failure (empirical).
    \item Soft-Exp inference preserves discrete support (Lemma 5.5), avoiding drift.
    \item Soft-Exp therefore works (empirical).
\end{itemize}

\subsection{What a Truly Closed Architecture Requires}

Theorem 5.3 shows that patchwork is impossible. A truly closed architecture requires:
\begin{enumerate}
    \item Continuous input (no discrete tokens)
    \item Closed internal layers
    \item Continuous output (no discrete token output)
\end{enumerate}

This is a significant departure from the autoregressive Transformer. It is the subject of the CWF Manifesto research charter.

% =============================================================================
\section{CWF Phase 1: The Decisive Experiment}
% =============================================================================

\subsection{Why Lorenz?}

The Lorenz system
\begin{align*}
    \dot{x} &= \sigma(y - x) \\
    \dot{y} &= x(\rho - z) - y \\
    \dot{z} &= xy - \beta z
\end{align*}
with $\sigma = 10$, $\rho = 28$, $\beta = 8/3$, has maximum Lyapunov exponent $\lambda \approx 0.906$. For a model to ``learn the Lorenz system'', it must either match the true ODE dynamics or learn an internal representation with the same Lyapunov structure.

The Oracle baseline (RK4 integration) provides the physics upper bound: EPT@0.9 = 100 steps.

\subsection{The Three Models}

\textbf{CWF}: A purpose-built continuous architecture with the six closed components. 0.98M parameters.

\textbf{AR+VQ}: Standard autoregressive Transformer with a VQ-VAE bottleneck (codebook size 512) on the input. 0.76M parameters.

\textbf{Oracle}: Direct RK4 integration. No parameters, no training.

\subsection{Training Protocol}

All models trained for 500 steps on 200 trajectories (1024 steps each) generated by RK4. Sequence length 256. AdamW, LR $1 \times 10^{-3}$, cosine schedule.

\subsection{Results}

\begin{table}[h]
\centering
\small
\begin{tabular}{lcccccc}
\toprule
Model & MSE@1 & MSE@10 & MSE@50 & MSE@100 & EPT@0.9 & Efficiency \\
\midrule
CWF & $24.5$ & $24.8$ & $57.6$ & $46.2$ & $\mathbf{2}$ & $0.02$ \\
AR+VQ & $224$ & $282$ & $257$ & $253$ & $\mathbf{3}$ & $0.03$ \\
Oracle & $0.18$ & $0.29$ & $0.19$ & $0.20$ & $\mathbf{100}$ & $1.00$ \\
\bottomrule
\end{tabular}
\caption{100-step rollout MSE and EPT for three models on the Lorenz system. CWF EPT@0.9 = 2 steps, AR+VQ = 3, Oracle = 100. Neither learned model is in-distribution for chaotic dynamics.}
\label{tab:lorenz}
\end{table}

By the pre-registered GENUINE\_PASS criterion (EPT efficiency $\geq 0.5$), CWF fails by a factor of 25.

\subsection{Diagnostic: Why CWF Failed on Rollout}

The rollout failure, despite one-step prediction success, indicates that CWF's forward pass is not learning the \emph{dynamics} of the Lorenz system. The architecture lacks:

\begin{itemize}
    \item Time-step embedding (no $\Delta t$ input)
    \item Recurrent state propagation (no $h_{t+1} = f(h_t, x_t)$)
    \item ODE solver (no numerical integration of learned dynamics)
\end{itemize}

Without these, CWF is a ``sequence mapper'' that learns to predict the next state from a context window, not a ``dynamics learner'' that models continuous evolution.

\subsection{The Verdict}

CWF achieved partial success: closure was verified, one-step prediction worked, architecture was implemented correctly. But CWF failed the GO/NO-GO gate: EPT@0.9 was only 2 steps, far from Oracle's 100 steps.

The failure was \textbf{not} in the closure property itself; closure was maintained throughout training. The failure was in the absence of ODE evolution machinery.

\subsection{What Would Save CWF}

To pass the GO/NO-GO gate, CWF would need:
\begin{enumerate}
    \item Neural ODE solver
    \item Time-step embedding
    \item Recurrent state propagation
\end{enumerate}

These additions are non-trivial engineering (1--2 person-years).

\subsection{What This Proves and Does Not}

\textbf{What this proves}: A closed continuous architecture (verified closure) does not automatically solve continuous dynamics tasks. Closure is necessary but not sufficient.

\textbf{What this does not prove}: The CWF Manifesto hypothesis is not falsified; only the minimal CWF prototype (closure without ODE evolution) is.

% =============================================================================
\section{Where Continuity Survived: The Soft-Exp Survivor}
% =============================================================================

\subsection{What Soft-Exp Is}

Soft-Exp is a one-line modification to autoregressive inference:
$$\mathbf{e}_{\text{next}} = \mathbf{p}^\top \mathbf{E}$$
where $\mathbf{p}$ is the softmax probability vector and $\mathbf{E}$ is the embedding matrix.

\subsection{The Empirical Result}

On V49 1.2B: argmax PPL $64.74$, Soft-Exp PPL $33.27$, $+48.6\%$ improvement. Scale-invariant across 2M--1.2B.

\subsection{Why Soft-Exp Avoids the Failure Mode}

By Lemma 5.5, the Soft-Exp expected embedding lies in the convex hull of the discrete embedding image. So $\mathrm{supp}(p_{\text{infer}}) \subseteq \mathrm{supp}(p_{\text{train}})$. The inference distribution is narrower than (not broader than) the training distribution, avoiding drift.

\subsection{Why Soft-Exp Training Failed}

Exp 30 tested training-time Soft-Exp: catastrophic failure. Training on a continuous input distribution causes overfitting (Corollary 5.4).

\subsection{The Unique Position of Soft-Exp}

Soft-Exp combines:
\begin{enumerate}
    \item Continuous feedback
    \item No training modification
    \item No retraining
    \item Scale-invariant improvement
\end{enumerate}

These properties follow from the closure framework. Soft-Exp respects the autoregressive framework while exploiting information-theoretic opportunity.

% =============================================================================
\section{Conclusion and Research Challenge}
% =============================================================================

\subsection{Summary of 35 Experiments}

The single positive result is Soft-Exp inference. The 25 failures span the entire patching design space. The closure argument (§5) explains why: patching cannot be closed, and training-time continuous feedback causes drift.

\subsection{The Three Lessons}

\begin{enumerate}
    \item Patching is structurally futile (closure violation).
    \item Closure is necessary but not sufficient (ODE evolution also required).
    \item Simple interventions often beat architectural redesign (Soft-Exp is one line of code).
\end{enumerate}

\subsection{The CWF Manifesto as Research Challenge}

A truly continuous language model would require:
\begin{enumerate}
    \item Continuous input (no discrete tokenization)
    \item Continuous internal state in $\mathbb{D}^d$
    \item Closed layers
    \item ODE evolution
    \item Time-step embedding
    \item Recurrent state
    \item Continuous output
\end{enumerate}

Building this is a 1--2 person-year investment. The CWF Manifesto remains an open research charter.

\subsection{What We Recommend}

\begin{enumerate}
    \item Do not patch continuous components into autoregressive Transformers.
    \item Do consider Soft-Exp inference for production deployments.
    \item Do investigate fully continuous architectures, with realistic expectations.
    \item Do share negative results.
    \item Do not assume mathematical elegance implies engineering viability.
\end{enumerate}

\subsection{The Companion Paper}

This paper is the long-form record. A companion paper documents the Soft-Exp method in detail \citep{wang2026softexp}. Together, the two papers tell a complete story.

\subsection{Final Acknowledgments}

This work represents three months of CMT experimentation plus two weeks of CWF experimentation. The 30 CMT rounds were not wasted; each informed the next. The 5 CWF rounds were not futile; they falsified a hypothesis. The single Soft-Exp survivor is what survived 34 rounds of systematic elimination.

% =============================================================================
\section*{Acknowledgments}
% =============================================================================

We thank the open-source community for PyTorch, NumPy, and SciPy. We thank the reviewers who pushed us to formalize the closure argument and design the Phase 1 decisive experiment.

% =============================================================================
% References
% =============================================================================

\begin{thebibliography}{99}

\bibitem{nitta2003complex}
T. Nitta.
\textit{On the critical points of the complex-valued neural network}.
ICONIP, 2003.

\bibitem{trabelsi2018deep}
C. Trabelsi et al.
\textit{Deep complex networks}.
ICLR, 2018.

\bibitem{liu2024kan}
Z. Liu et al.
\textit{KAN: Kolmogorov-Arnold networks}.
arXiv preprint, 2024.

\bibitem{lezama2018overcoming}
J. Lezama.
\textit{Overcoming the disentanglement problem: Self-supervised visual content}.
CVPR, 2018.

\bibitem{wang2023rope}
S. Wang et al.
\textit{RoFormer: Enhanced transformer with rotary position embedding}.
Neurocomputing, 2023.

\bibitem{born1926quantenmechanik}
M. Born.
\textit{Zur Quantenmechanik der Sto\ss{}vorg\"ange}.
Zeitschrift f\"ur Physik, 1926.

\bibitem{bengio2015scheduled}
S. Bengio, O. Vinyals, N. Jaitly, N. Shazeer.
\textit{Scheduled sampling for sequence prediction with recurrent neural networks}.
NeurIPS, 2015.

\bibitem{hinton2015distill}
G. Hinton, O. Vinyals, J. Dean.
\textit{Distilling the knowledge in a neural network}.
NIPS Deep Learning Workshop, 2015.

\bibitem{huszar2015how}
F. Huszár.
\textit{How (not) to train your generative model: Scheduled sampling, likelihood, adversary?}.
arXiv:1511.05151, 2015.

\bibitem{chen2018neuralode}
R. T. Q. Chen, Y. Rubanova, J. Bettencourt, D. Duvenaud.
\textit{Neural ordinary differential equations}.
NeurIPS, 2018.

\bibitem{wang2026softexp}
Y. Wang.
\textit{Exposure bias is bigger than you think: Continuous expected embeddings as a drop-in fix for autoregressive decoding}.
CrystaLLM Project, 2026.

\bibitem{wang2026cwf}
Y. Wang.
\textit{Closed Waveformer: A research manifesto}.
CrystaLLM Project, 2026.

\bibitem{v49scale}
Y. Wang.
\textit{V49 scale: A 1.2B-parameter Transformer baseline for autoregressive language modeling}.
CrystaLLM Project, 2026.

\end{thebibliography}

% =============================================================================
% Appendix A: Cross-Cutting Failure Modes
% =============================================================================

\appendix

\section{Cross-Cutting Failure Modes}

\subsection{Failure Mode 1: Training-Inference Distribution Drift}

When the training input distribution differs from the inference input distribution, the model overfits to training. At inference, it sees OOD inputs and either memorizes or fails catastrophically.

\textbf{Empirical evidence}: CMT Exp 7--25 (memorization at 8--12k steps), Exp 30 (Soft-Exp training, catastrophic). \textbf{Root cause}: distribution shift within the model's input manifold.

\subsection{Failure Mode 2: Closure Violation in Patchwork Architectures}

A patchwork architecture combines discrete components (token embedding, vocab head) with continuous components (complex attention, KAN FFN) without ensuring the entire pipeline maintains closure. The interface creates information bottlenecks and gradient instability.

\textbf{Empirical evidence}: All CMT Exp 6--25. \textbf{Root cause}: discrete-continuous interface leaks through every continuous component (Lemma 5.2).

\subsection{Failure Mode 3: Architectural Absence of ODE Evolution}

The model lacks Neural ODE solver, recurrent state, time-step embedding. Without these, it can only do discrete ``step-to-step'' prediction, which fails on chaotic dynamics tasks.

\textbf{Empirical evidence}: CWF Phase 1 (Lorenz, EPT = 2 vs Oracle 100). \textbf{Root cause}: single-shot forward pass, no continuous evolution.

\subsection{Failure Mode 4: Tokenization Bottleneck on Continuous Data}

VQ bottleneck on continuous data: $V$ codes cannot represent a 3D continuous chaotic orbit. Information loss is unbounded.

\textbf{Empirical evidence}: AR+VQ on Lorenz (MSE stuck at 232). \textbf{Root cause}: information-theoretic bound on quantization.

\subsection{Failure Mode 5: Optimizer-Landscape Mismatch on Complex Losses}

Real AdamW on complex-Wirtinger gradients treats 2D gradient as 1D, losing rotational structure. Oscillation, sharp minima, memorization.

\textbf{Empirical evidence}: CMT phase transition at 10--12k steps. \textbf{Root cause}: standard optimizer insufficient for complex-valued models.

% =============================================================================
% Appendix B: Reproducibility
% =============================================================================

\section{Reproducibility}

\subsection{Code Repositories}

\begin{itemize}
    \item \texttt{experiments/v49\_pre/}: CMT and Soft-Exp experiments (exp01--exp30)
    \item \texttt{research/cwf/}: CWF Manifesto and experiments
    \item \texttt{research/cwf/experiments/exp02\_lorenz/}: Lorenz Phase 1
    \item \texttt{docs/papers/}: Paper drafts and figures
\end{itemize}

\subsection{Data}

\begin{itemize}
    \item Training: \texttt{bpe\_train\_65M\_s42.npy} (65M BPE tokens)
    \item Validation: \texttt{bpe\_val\_s42.npy}
    \item Character-level: \texttt{v28\_train.parquet}, \texttt{v28\_val.parquet}
    \item CWF Phase 1: RK4-generated Lorenz trajectories on GPU
\end{itemize}

\subsection{Hardware}

All experiments on RTX 5090 (32 GB). Peak memory 4 GB for CWF, 2 GB for others.

\subsection{Hyperparameters}

See per-experiment details in supplementary material. All seeds fixed at 42.

\subsection{Compute Cost}

Total: $\sim$350 GPU-hours ($\sim$1 GPU-month).

\subsection{Software}

Python 3.10, PyTorch 2.9.1+cu128, NumPy 2.2.6, SciPy 1.13, tiktoken 0.7, matplotlib 3.10.8.

% =============================================================================
% Appendix C: Timeline Table
% =============================================================================

\section{Timeline of 35 Experiments}

\subsection{CMT Phase (Exp 1--25)}

\begin{table}[h]
\centering
\small
\begin{tabular}{p{1.5cm}p{3.5cm}p{3.5cm}p{3cm}p{1.5cm}}
\toprule
ID & Hypothesis & Setup & Key Result & Verdict \\
\midrule
1 & Mamba3/SSD replaces attention & Selective SSM & PPL higher & INCONCLUSIVE \\
2 & Complex KAN beats MLP & KAN complex & Worse than GELU & FAIL \\
3 & FP8 mixed precision & FP8 forward & Speedup, no quality & NEUTRAL \\
4 & 8-bit AdamW & 8-bit states & No regression & NEUTRAL \\
5 & Curriculum learning & Easy-to-hard & No improvement & INCONCLUSIVE \\
6 & CMT-FFN only & KAN complex & PPL 4.5 & FAIL \\
7 & CMT-full sanity & All 3 knives & Memorizer @ 4k & FAIL \\
8 & CMT-full refined & Bug fixes & Memorizer @ 8k & FAIL \\
9 & Baseline rerun & Sanity & Stable & BASELINE OK \\
10 & CMT softmax fix & Fix softmax & Marginal & FAIL \\
11 & CMT KAN complex mul & Fix multiply & No improvement & FAIL \\
12 & CMT LieRE Cayley & Real Cayley & No help & FAIL \\
13 & CMT real init v2 & Better init & Marginal & FAIL \\
14 & CMT no context PE & Skip ctx PE & PPL 32.58 peak & FAIL \\
15 & CMT-full v2 & All fixes & Memorizer @ 12k & FAIL \\
16 & CMT-clean & 0-bug & Memorizer @ 30k & FAIL \\
17 & Phase transition & 4 ckpts & Transition at 10k & INSIGHT \\
18 & A1 long training & 8k LR 3e-5 & Underfit & FAIL \\
19 & 5k sanity & 2M 2k samples & Loss 33$\to$18 & PARTIAL \\
20 & BPE sanity & BPE vocab 4100 & Coherent 0/6$\to$2/6 & MILD \\
21 & Baseline BPE 5k & Non-CMT & Coherent 1/6 & BASELINE \\
22 & CMT BPE 16k & Long training & PPL rebound & MISLEADING \\
23 & CMT BPE 10k+16k & More data & Rebound 0.00x & CORRECTION \\
24 & Cayley PE isolated & vs RoPE/NoPE & Not better & PARTIAL \\
25 & CMT BPE 10k+32k & Final CMT & Memorizer & CMT\_DEAD \\
\midrule
\end{tabular}
\end{table}

\subsection{Soft-Exp Phase (Exp 26--30)}

\begin{table}[h]
\centering
\small
\begin{tabular}{p{1.5cm}p{3.5cm}p{3.5cm}p{3cm}p{1.5cm}}
\toprule
ID & Hypothesis & Setup & Key Result & Verdict \\
\midrule
26 & Soft-Exp 2M & Inference & $+81\%$ & PASS \\
27 & Soft-Exp 8M & Boundary scan & $+53\%$ & PASS \\
28 & Soft-Exp 16M/32M & LR fix & 16M $+48.6\%$ & PASS \\
29 & Soft-Exp V49 1.2B & Real scale & $+48.6\%$ & STRONG \\
30 & Soft-Exp at training & $\alpha = 0.5$ & Catastrophic & FAIL \\
\bottomrule
\end{tabular}
\end{table}

\subsection{CWF Phase (Exp 01--02)}

\begin{table}[h]
\centering
\small
\begin{tabular}{p{2cm}p{3.5cm}p{3.5cm}p{3cm}p{1.5cm}}
\toprule
ID & Hypothesis & Setup & Key Result & Verdict \\
\midrule
CWF 01 & Harmonic prediction & CWF vs AR & Converges faster & MIXED \\
CWF 02.0 & Lorenz data + Oracle & RK4 & Oracle EPT=100 & ORACLE OK \\
CWF 02.1 & 3-channel CWF & Adapter & Closure OK & PASS \\
CWF 02.2 & AR+VQ baseline & VQ bottleneck & Codebook collapse & FAIL \\
CWF 02.3 & 100-step rollout & 6 metrics & CWF EPT=2 & TOTAL\_FAIL \\
\bottomrule
\end{tabular}
\end{table}

\subsection{Summary}

\begin{table}[h]
\centering
\small
\begin{tabular}{lc}
\toprule
Verdict & Count \\
\midrule
STRONG PASS & 1 \\
PASS & 4 \\
PARTIAL & 1 \\
INCONCLUSIVE & 4 \\
FAIL & 25 \\
\midrule
\textbf{Total} & \textbf{35} \\
\bottomrule
\end{tabular}
\caption{Verdict distribution across 35 experiments.}
\end{table}

\end{document}
